\begin{example}
    La sucesión $\partfrac{\sin n}$ es densa en $[0,1)$ pero no es equidistribuida.
\end{example} 

\begin{resolve}
    Sabemos que $n$ es equidistribuida en $[0,2\pi)$.
    Por la caracterización funcional de la equidistribución aplicada a la función $f(x) = e^{2\pi i \sin x}$, tenemos que
    \[
        \frac{1}{N} \sum_{n=1}^N e^{2\pi i \sin n} \longrightarrow \frac{1}{2\pi} \int_0^{2\pi} e^{2\pi i \sin x} dx \qquad \text{cuando } N \to \infty.
    \]

    Si calculamos esta integral, tenemos que
    \[
        \int_0^{2\pi} e^{2\pi i \sin x} dx = 2\pi J_0(2\pi) \neq 0.
    \]

    Si $\partfrac{\sin n}$ fuese equidistribuida, por el criterio de Weyl, debería cumplirse
    \[
        \forall \, k \in \mathbb{Z}\setminus\{0\}, \quad \frac{1}{N} \sum_{n=1}^N e^{2\pi i k \sin n} \longrightarrow 0 \qquad \text{cuando } N \to \infty,
    \]
    Pero para $k=1$ hemos visto que esta suma tiende a $J_0(2\pi) \neq 0$, lo que demuestra que $\partfrac{\sin n}$ no es equidistribuida.
\end{resolve}